\subsubsection{Agenten und KI-Agenten}
Zum vollen Verständnis ist es wichtig, zwischen den verschiedenen Begriffen der Agenten zu unterscheiden. Agenten im Allgemeinen sind autonome Systeme, die fortlaufend Informationen aus ihrer Umgebung aufnehmen, diese intern verarbeiten und darauf basierend zielgerichtete Handlungen ausführen. Die grundlegende Struktur dafür besteht aus zwei Komponenten: den Sensoren, über die Informationen aus der Umwelt aufgenommen werden, und den Aktuatoren, über die der Agent Aktionen ausführt. Dadurch entsteht ein kontinuierlicher Kreislauf aus Wahrnehmung und Aktion, in dem der Agent auf Basis der aufgenommenen Informationen Entscheidungen trifft und durch seine Handlungen aktiv auf die Umgebung einwirkt. \cite[\pagef 34f]{russell_artificial_2016}

Darüber hinaus ist es wichtig zu erwähnen, dass ein Agent aus einer Agentenfunktion, einem Agentenprogramm und der zugrunde liegenden Architektur besteht. Die Agentenfunktion beschreibt auf abstrakter Ebene, wie ein Agent aus seinen Wahrnehmungen eine Entscheidung ableitet und eine geeignete Handlung auswählt. Das Agentenprogramm stellt die konkrete Implementierung dieser Funktion dar und wird in Form einer Software realisiert, die den Entscheidungsprozess steuert und die Ausführung der Handlungen übernimmt. Die zugrunde liegende Architektur stellt die notwendigen technischen Mittel bereit, um Wahrnehmungen zu erfassen, zu verarbeiten und Handlungen auszuführen. Sie umfasst sowohl die sensorischen und aktorischen Schnittstellen zur Umwelt als auch die Rechenressourcen, auf denen das Agentenprogramm betrieben wird. Erst durch das Zusammenspiel dieser drei Komponenten entsteht ein vollständig funktionsfähiger Agent. \cite[\pagef 35]{russell_artificial_2016}

Zusammenfassend lässt sich ein Agent als ein System beschreiben, das Informationen aus seiner Umgebung verarbeitet und daraus geeignete Handlungen ableitet. Dieses grundlegende Funktionsprinzip bildet die Basis für unterschiedliche Ausprägungen von Agenten, die sich in ihren Eigenschaften unterscheiden. Vor diesem Hintergrund können Agenten in verschiedene Kategorien eingeteilt und entsprechend klassifiziert werden. \cite[\pagef 22]{ertel_grundkurs_2024}

Auf Grundlage des allgemeinen Agentenmodells lassen sich Software-Agenten als eine zentrale Kategorie einordnen. Nwana (1996) beschreibt sie als autonome Softwareeinheiten, die im Auftrag eines Nutzers oder eines anderen Programms Aufgaben ausführen und dabei vollständig in digitalen Umgebungen agieren. Sie verfügen über keine materiellen Sensoren oder Aktuatoren, sondern ihre Wahrnehmungs- und Handlungsmöglichkeiten ergeben sich ausschließlich aus digitalen Schnittstellen. Die zugrunde liegende Architektur besteht folglich aus den Rechen- und Kommunikationsressourcen, auf denen das Agentenprogramm ausgeführt wird. \cite[\pagef 205ff]{nwana_software_1996}

Agenten im Allgemeinen, aber auch Software-Agenten, lassen sich zusätzlich in verschiedene Unterkategorien einteilen, die nicht strikt voneinander getrennt sind, sondern vielmehr unterschiedliche Eigenschaften und Fähigkeiten von Agenten beschreiben. Als erste Kategorie lassen sich Reflex-Agenten identifizieren, die lediglich auf die aktuell vorliegende Eingabe reagieren und keine Gedächtnisfunktion besitzen. Im Gegensatz dazu stehen die modellbasierten Agenten oder auch Agenten mit Gedächtnis genannt, welche zusätzlich die Vergangenheit mit berücksichtigen können. Darüber hinaus existieren zielorientierte Agenten, die explizite Ziele verfolgen und ihre Handlungen entsprechend ausrichten. Darauf aufbauend werden nutzenorientierte Agenten unterschieden, die zusätzlich versuchen, den langfristigen Nutzen ihrer Entscheidungen zu maximieren. Darüber hinaus sind noch die lernfähigen Agenten wichtig zu nennen. Diese Agenten sind in der Lage, sich selbst durch erfolgreiche Aktionen oder durch Feedback anzupassen und somit den Nutzen ihrer Aktionen im Laufe der Zeit zu verbessern. \cite[\pagef 23f]{ertel_grundkurs_2024}

Die für diese Arbeit im Mittelpunkt stehende Kategorie an Agenten sind die \ac{KI}-Agenten. Die Definition von diesen ist, genauso wie bei der Künstlichen Intelligenz, nicht einheitlich. Wie bereits erwähnt, handelt es sich bei Agenten um Systeme, die ihre Umgebung wahrnehmen und auf Basis dieser Wahrnehmungen selbstständig Handlungen ausführen, um festgelegte Ziele zu erreichen \cite[\pagef 4]{russell_artificial_2016}. Der Begriff der \ac{KI}-Agenten wird in moderner Literatur häufig um den Begriff der Sprachmodelle erweitert. So beschreiben Koch et al. (2025), dass \glqq{} [e]in \ac{KI}-Agent im Kern ein Sprachmodell (LLM) [ist], das mit der Möglichkeit ausgestattet wurde, relevante Informationen zu suchen, zu speichern (Memory) und mit der Außenwelt (andere Systeme oder KI-Agenten) interagieren kann.\grqq{} \cite[\pagef 111]{koch_ki-agenten_2025}
Dadurch, dass die \ac{LLM} mit der Fähigkeit zur Informationssuche, Speicherung und Interaktion mit der Außenwelt ausgestattet sind, können \ac{KI}-Agenten komplexe Aufgaben und Probleme lösen, die typischerweise menschliche Intelligenz erfordern. Darüber hinaus ermöglicht das \ac{LLM} eines \ac{KI}-Agenten im Unterschied zu anderen Software-Agenten ein autonomes Handeln und verleiht dem Agenten eine höhere Flexibilität und Anpassungsfähigkeit. Sie unterstützen ein lernbasiertes Verhalten und erlauben die kontinuierliche Verfeinerung von Aktionen auf Grundlage von Feedback. \cite[\pagef 13]{karrenberg_technologische_2025} 
Die große Besonderheit liegt also genau in den gerade genannten Punkten von \ac{KI}-Agenten. Anders als andere Agenten sind diese in der Lage, komplexe Probleme zu lösen, ohne dass jeder einzelne Schritt explizit vorgegeben werden muss. \cite[\pagef 3]{karrenberg_technologische_2025} Zusammenfassend lässt sich festhalten, dass Large Language Models in modernen \ac{KI}-Agenten eine zentrale Rolle einnehmen, jedoch kein eigenständiges Agentensystem darstellen. Vielmehr fungiert ein LLM als kognitive Kernkomponente, die für das Verstehen, Planen und Ableiten von Entscheidungen verantwortlich ist. Wahrnehmung und Handlung werden hingegen durch zusätzliche Systemkomponenten realisiert, etwa durch Schnittstellen zur Informationsbeschaffung oder zur Ausführung von Aktionen in externen Systemen. Xi et al. (2025) beschreiben dieses Zusammenspiel, indem sie das LLM als Bestandteil eines sogenannten brain module einordnen, während Wahrnehmung und Handlung durch separate Module umgesetzt werden. Erst durch die Kombination dieser Komponenten entsteht ein vollständig funktionsfähiger \ac{KI}-Agent im Sinne der klassischen Agentendefinition. \cite[\pagef 1ff]{xi_rise_2025}
